\section{Return maps and the cubic comparison}\label{sec:foundations}
We use the right-half-plane fields $X_s=(y-F(sx),-x)$, $s=\pm1$.
Time runs forward for $s=+$ and backward in the original left half-plane
for $s=-$. Write $\Sigma_-$ for the open negative $y$-axis. The following forms of the return estimates specify the
inputs needed below. Their proofs are in Sections~2 and~3 of the main article~\cite{MainArticle}.

\subsection{Axis hits and a common linear parameter}
\begin{lemma}[Half-returns]\label{lem:half-returns}
(i) Each periodic orbit meets the open positive $y$-axis in exactly one point.
(ii) Let $B_s(q,h)>0$ be the magnitude of a first lower-axis hit from $(q,h)$,
and write $w_s(h)=B_s(0,h)$. At transverse first hits these functions
are analytic. (iii) Between two returning axis launches the whole strip
returns, with $w_s'(h)>0$. If $|F'(0)|<2$, one upper trap suffices down
to zero. If also $F(sx)\to+\infty$, every axis launch $h>0$ and every
off-axis launch $q>0$ has a lower hit. At fixed ordinate these hits are
strictly ordered in $q$ on their domain.
\end{lemma}
Indeed, $y$ decreases strictly in $x>0$. Returning arcs are disjoint
graphs there, with $\dot x=y\ne0$ at each axis crossing. Their order
gives the trapping strip, and the angular velocity near the origin is
strictly negative when $|F'(0)|<2$. At large positive $x$, coercivity
bounds every forward arc. In the interior the variables $u=x^2/2$ and
$x$ give
\begin{align}
 \frac{\dd u}{\dd(-y)}&=y-F(s\sqrt{2u}),\label{eq:scalar-u}\\
 \frac{\dd y}{\dd x}&=\frac{x}{F(sx)-y}\quad\text{on an inward tail}.\label{eq:spatial-ode}
\end{align}

\begin{lemma}[Incoming trajectories]\label{lem:incoming-rays}
For even-degree $P$, $m=\deg P\ge4$, $P(0)=0$, with positive leading
coefficient and $|P'(0)|<2$, an incoming trajectory from infinity at
ordinate $y$ has a unique transverse first hit $-I_{s,P}(y)$.
(i) Its hit varies continuously with the ordinate and coefficients at
fixed degree. (ii) If $P_\nu\to P$ in $C^1$ on compact sets and
$P_\nu(sx)\ge c x^m$ on $[X_0,Z_\nu]$, $Z_\nu\to\infty$, launches
from $(Z_\nu,y)$ converge in hit magnitude uniformly on compact ordinate
intervals. (iii) $I_{s,P}(-w)=w+O(w^{2/m-1})$ as $w\to\infty$.
\end{lemma}
Inward integration of~\eqref{eq:spatial-ode} contracts vertical
differences. Above a large fixed $X$ the drift is $O(X^{2-m})$
uniformly on a bounded ordinate interval. Below $X$ one uses compact
flow dependence up to a transverse hit. The same proof applies near
any fixed finite list of transverse hits for the limiting polynomial,
without a condition on its slope at zero (main paper, Corollary~2.2a).

\begin{lemma}[Sensitivity of a half-return]\label{lem:rotational}
For $F_a(x)=F_0(x)+ax$, $G_{s,a}(x)=F_a(sx)$, and a transverse first
hit at time $T$ from a fixed launch,
\begin{equation}\label{eq:rot-half}
 B_s\,\partial_a B_s=-s\int_0^T x(t)^2
          \exp\!\left(-\int_t^T G_{s,a}'(x(v))\,\dd v\right)\dd t .
\end{equation}
Thus $\partial_a\D<0$ wherever both axis half-returns exist.
\end{lemma}
This follows by varying the flow and projecting to the axis at the hit.
The variational determinant is $-s x^2$ and the divergence is $-G_{s,a}'$.

\begin{proposition}[Simultaneous reference detuning, Sard regularization, and odd-degree barrier]\label{prop:sard}
Let $F_{a_*}(x)=\sum_{j=2}^n a_j x^j + a_* x$ have exact degree $n\ge 2$ ($a_n\neq 0$).
\begin{enumerate}
  \item[\textup{(i)}] Suppose $0<h_0<h_1<\dots<h_N$ lie in an open interval $I_h\subset(0,\infty)$ on which both transverse half-returns $w_\pm(h,a_*)$ exist. Assume there is a common sign $\sigma_{\mathrm{det}}\in\{+1,-1\}$ such that, after setting $\sigma_i=\sgn\D(h_i,a_*)$ when $\D(h_i,a_*)\neq 0$ and $\sigma_i=\sigma_{\mathrm{det}}$ at every matched reference height where $\D(h_i,a_*)=0$, the sequence $(\sigma_0,\dots,\sigma_N)$ strictly alternates ($\sigma_{i-1}\sigma_i=-1$ for $1\le i\le N$). Suppose in addition that $F_{a_*}$ already possesses $q_0\ge 0$ hyperbolic limit cycles contained in a compact region $K_0\subset\R^2$ disjoint from $\{(0,h):h\in[h_0,h_N]\}$. Then there is a regular value $a_{\rm reg}$ arbitrarily close to $a_*$ and an open neighborhood $\mathcal{V}\subset\{(c_1,\dots,c_n)\in\R^n:c_n\neq 0\}$ of the resulting field on which every system~\eqref{eq:lienard} has at least $q_0+N$ geometrically distinct hyperbolic limit cycles.
  \item[\textup{(ii)}] Suppose $|a_*|<2$, $n$ is even, and $0<h_0<\dots<h_N$ have strict alternating signs $\D(h_i,a_*)$ ending with $\D(h_N,a_*)>0$. Then for each sufficiently small $\rho>0$ one can choose $a_{\mathrm{reg}}$ arbitrarily close to $a_*$ such that every system in an open exact-degree-$(n+1)$ neighborhood of $F_{a_{\mathrm{reg}}}(x)+\rho x^{n+1}$ has at least $N+1$ distinct hyperbolic limit cycles.
\end{enumerate}
Each open coefficient neighborhood in \textup{(i)} or \textup{(ii)} contains a rational vector by density.
\end{proposition}
Here $\D=w_+-w_-$. On the zero set $\partial_a\D<0$ gives
a local graph $a=\alpha(h)$. A single regular value on the trapped
return intervals makes the zeros simple, since $w_-'(h)>0$.
Part (ii) uses the outer barrier in the main paper's Proposition~2.4.

\subsection{Actions and exterior transfer}
Set $\phi(u)=\Phi(\sqrt{2u})$ for a scaled half-profile $\Phi(z)$.
On a rising branch shared with the reference, the \emph{action} is the
magnitude of the reduced slow divergence integral. Its density in the
increasing spatial coordinate is $\Phi'(z)^2\,\dd z/z=\phi'(u)^2\,\dd u$. An incoming trajectory follows the
outermost rising root. At a minimum it jumps to the transverse inner
landing, deleting the intervening open spatial interval. An axis launch
inside a fold window meets the reference on the inner branch below that
landing. With $r$ the landing of the window containing $Y$, this gives
\begin{equation}\label{eq:JU-def}
 J^U(Y)=
 \begin{cases}
  J^I(Y),&Y\ \text{outside the fold windows},\\
  \text{retained action up to }r,&Y\ \text{inside a fold window}.
 \end{cases}
\end{equation}
For two ordered solutions of $\epsilon u_t=\phi(u)-t$ one has exactly
\begin{equation}\label{eq:secant-exact}
 \epsilon(\log(u_1-u_2))'=[\phi]_{u_1,u_2}
     =\frac{\phi(u_1)-\phi(u_2)}{u_1-u_2}.
\end{equation}

\begin{proposition}[Matched-maximum exterior transfer]\label{prop:cap-transfer}
Let $P\in\R[x]$ have $P(0)=0$, even degree $m=\deg P\ge 6$ with leading coefficient $c=[x^m]P>0$, and highest odd power $j=m-g$ with $g\in\{1,3\}$ and $b=[x^j]P\neq 0$. Suppose $P$ possesses $q\ge 0$ hyperbolic limit cycles, and at a nonempty strictly decreasing finite list of ordinates $y_0>y_1>\dots>y_r$ and a real shift $\Delta\in\R$, both incoming trajectories in $I_{+,P}(y_i)-I_{-,P}(y_i-\Delta)$ exist, hit $\Sigma_-$ transversely, and exhibit $r$ strict sign changes ending with $\sigma=\sgn\big(I_{+,P}(y_r)-I_{-,P}(y_r-\Delta)\big)\in\{\pm 1\}$. Put $p=2-g\in\{+1,-1\}$. Then there is a nonempty open subset of exact-degree-$\max(m+2,j+4)$ polynomials (including rational polynomials) on which every field~\eqref{eq:lienard} has at least
\begin{equation}\label{eq:cap-count}
  q + r + 1 + \ind_{\{\sigma\neq\sgn(bp)\}}
\end{equation}
simultaneous hyperbolic limit cycles, retaining all $q$ compact cycles while placing the new upper crossings near $R^m\frac{2c}{m+2}\to+\infty$.
\end{proposition}
The two exterior coefficients match the pair of reference hits
and one entrance pair. On each side the remaining entrances share
the same terminal map $K$. At an intermediate regular height $b$
their separations from the reference satisfy
\begin{equation}\label{eq:entry-distortion}
 \epsilon\log d(\alpha)=\int_\alpha^b\phi'(u(t))\,\dd t
                    +O(\epsilon|\log\epsilon|),
\end{equation}
on the attracting passage in increasing $t$, with $u(t)$ the falling
root. Differentiation of the upper-axis first hit gives
\begin{equation}\label{eq:Ks-deriv}
 K'(\xi)=\frac{\epsilon}{K(\xi)}\exp\!\left(
 \epsilon^{-1}\int_b^{K(\xi)}\phi'(U(t;\xi))\,\dd t\right).
\end{equation}

\begin{lemma}[Common terminal section, main paper, Lemma~3.4]\label{lem:terminal-distortion}
Let a whole section interval $B$ flow under $\epsilon U'=f(t,U)$ to a
first downward crossing of $U=u_c$ and then to zero. Before crossing,
assume bounded elapsed time and uniform bounds on $f,f_U,f_{UU}$;
in a fixed crossing collar assume $f\le-\nu<0$ and $f_t$ bounded.
On $[0,u_c]$ assume $f\le-\nu$ and $f,f_t,f_{tt}$ bounded and
continuous in the inverse chart. Take these bounds between the arcs,
with the regularity of the indicated derivatives on each piece of
the pre-crossing flow. If $D$ bounds the time integral of an envelope
of $(f_U)_+$ there and $|B|e^{D/\epsilon}$ is smaller than the
collar width, then $\operatorname{osc}_B(\epsilon\log K')\le
C|B|e^{D/\epsilon}$ and $K'>0$.
\end{lemma}
The estimate compares the two positive gaps throughout the interval
spanned by the reference and the launches. In the exterior construction
of the main paper the absolute gap expansion calibrates the two
coefficients, and the gap-ratio expansion resolves physical entrances
at scale $\tau^m$. Section~S5 gives the corresponding formulas for the
joint degree-twelve construction.

\subsection{Cubic data for the moving folds}
For $a=-3k^2<0$ ($k=\sqrt{-a/3}>0$) and $b\in\R$, the cubic $f(v)=v^3+av+b$ has a local maximum $M=b+2k^3$ at $v=-k$, a local minimum $\mu=b-2k^3$ at $v=k$, and rising-branch landing $v=-2k$ solving $f(-2k)=\mu$. Let $P_a(v)=\int_0^v (3r^2+a)^2\,\dd r=\tfrac95 v^5+2av^3+a^2v$, which is strictly increasing on $\R$ since $P_a'(v)=(f'(v))^2\ge 0$. For $\xi\notin\{\mu,M\}$, let $v(\xi)$ be the largest real root of $f(v)=\xi$, and define the incoming and upper cubic actions
\begin{equation}\label{eq:cubic-actions}
  J_{a,b}^I(\xi) = P_a(v(\xi)) - \ind_{\{\xi>\mu\}}\big(P_a(k)-P_a(-2k)\big),\qquad
  J_{a,b}^U(\xi) = \begin{cases}
    P_a(-2k), & \mu<\xi<M,\\
    J_{a,b}^I(\xi), & \xi\notin[\mu,M].
  \end{cases}
\end{equation}
Put $\Lambda=493/2$ and define the equal-weight cubic differences
\begin{equation}\label{eq:D-cub}
  D_{\mathrm{cub}}^X(\xi) = J_{-1,-8}^X(\xi) - J_{-3,8}^X(\xi) - \frac{493}{2},\qquad X\in\{U,I\}.
\end{equation}

Table~\ref{tab:cubic} gives the rational intervals used below. To
reproduce a row, isolate the largest root on its monotone branch by
the listed endpoints ($j/10^5,(j+1)/10^5$), then apply the increasing primitive $P_a$ to the two
endpoints. Inside an open fold window use $P_a(-2k)$ for the upper
launch. Lemma~B.1 of the main article gives the root and jump bounds.
\begin{table}[ht]
\centering\scriptsize\setlength{\tabcolsep}{3.5pt}
\caption{Certified open rational enclosures for $D_{\mathrm{cub}}^U(\xi)$ in~\eqref{eq:D-cub}. At $\xi\in\{-12,0,100,10^5\}$, $D_{\mathrm{cub}}^I(\xi)=D_{\mathrm{cub}}^U(\xi)$ satisfies the identical enclosure. The integers specify the root brackets in the proof; L denotes the inner landing.}\label{tab:cubic}
\smallskip
\begin{tabular}{@{}lrrrrrrrr@{}}
\toprule
$\xi$ & $-12$ & $-7.62$ & $-7.61$ & $0$ & $9.99$ & $10.01$ & $100$ & $10^5$ \\
\midrule
$j_+$ & $-179633$ & L & $115639$ & $216631$ & $274737$ & $274829$ & $483219$ & $4642430$ \\
$j_-$ & $-308086$ & $-289690$ & $-289645$ & $-249204$ & L & $200111$ & $473570$ & $4643619$ \\
Lower bound & $81$ & $-0.85$ & $0.40$ & $-79$ & $21$ & $-0.43$ & $610$ & $-98000$ \\
Upper bound & $82$ & $-0.82$ & $0.44$ & $-77$ & $23$ & $-0.40$ & $615$ & $-96000$ \\
Sign of $D_{\mathrm{cub}}^U(\xi)$ & $+$ & $-$ & $+$ & $-$ & $+$ & $-$ & $+$ & $-$ \\
Sign of $D_{\mathrm{cub}}^I(\xi)$ & $+$ &  &  & $-$ &  &  & $+$ & $-$ \\
\bottomrule
\end{tabular}
\end{table}

